\section{A research state with mathematical and historical structure}\label{sec:representation}

\subsection{Reasoning graphs and exploration histories}

An RRE state at revision $t$ is represented schematically as
\[
  \mathcal S_t=(\mathcal G_t,\mathcal H_t,\mathcal F_t,\mathcal A_t,\mathcal P_t),
\]
where $\mathcal G_t$ is the reasoning collection, $\mathcal H_t$ the exploration history, $\mathcal F_t$ the frontier of research targets, $\mathcal A_t$ the evidence and artifact store, and $\mathcal P_t$ the active objectives, operating mode, budget, and research and admission policies. Each objective records its origin, selection rationale, and acceptance conditions. In open-ended exploration, new objectives become explicit records before their investigation; subsequent revisions preserve their earlier versions.

The reasoning collection answers: \emph{What supports this mathematical statement?} The exploration history answers: \emph{What was attempted, when, using which available evidence, and with what outcome?} A proof may be reorganized after discovery into a shorter dependency graph. That reorganization should not rewrite the historical sequence of investigations. Conversely, an event that motivated a proof step is not necessarily a mathematical premise of that step.

The graph has typed nodes for definitions, assumptions, claims, constructions, counterexamples, obligations, and evidence records. An inference is a directed hyperedge
\begin{equation}\label{eq:hyperedge}
 e:\quad \Gamma_e;\ p_1,\ldots,p_k\ \Longrightarrow\ q,
\end{equation}
with all $p_i$ jointly required, an explicit scope $\Gamma_e$, and a linked justification. Alternative arguments for $q$ are separate hyperedges. The existence of an edge records an asserted inference, not automatic mathematical truth. Its status depends on the justification and on the statuses of its premises.

Scope includes domains, quantifiers, parameter dependence, and limiting regimes. A theorem for each fixed depth is different from a uniform theorem over growing depth. A result valid on a generic set is different from one valid at every parameter value. These distinctions are part of the mathematical content, not optional descriptive metadata. Where a common condition is inherited through an argument, a prerequisite trace must still expose it.

Navigation links record similarity, analogy, source relationships, or historical motivation. They never establish implication. This distinction is particularly important in an aggregated graph: two papers may share an instrument without one theorem supplying the hypotheses of the other. A visually short path is not necessarily a valid mathematical route.

\subsection{A glossary of status and review}

RRE separates four questions: what kind of mathematical assertion is recorded, what has been checked, how significant it might be, and whether it is new. A single ordered confidence scale cannot represent these distinctions adequately.

\begin{table}[htbp]
\centering\small
\begin{tabularx}{\linewidth}{@{}p{2.7cm}Y@{}}
\toprule
Mathematical status & Meaning in this paper \\
\midrule
Source assertion & A statement faithfully attributed to a source; independent validity is not implied. \\
Conjecture & A proposed statement without an admitted proof. \\
Conditional derivation & An argument with named assumptions or unresolved premises still attached. \\
Proof recorded & A complete argument is supplied at its declared level of rigor; its review record remains separate. \\
Computation & An exact or numerical result for stated inputs or a finite range. It supports only the claim actually checked. \\
Refutation & A counterexample or argument contradicts a specified statement under the same hypotheses. \\
Open obligation & A claim or quantitative bound is required by a route and has not been established. \\
\bottomrule
\end{tabularx}
\caption{Mathematical statuses. Supersession and withdrawal are revision states; they preserve the earlier record. A failed attempt is an event outcome, not a refutation.}
\label{tab:statuses}
\end{table}

Review is a collection of scoped records: source comparison, internal AI review, exact certificate replay, independent implementation, human review, or formal verification. These are not interchangeable. Exact arithmetic can certify a finite identity without establishing an all-size theorem. A formal kernel can validate an encoded theorem while leaving the correspondence to an informal source in need of review. Human review should identify who reviewed what, with which limitations; responsibility for a manuscript does not establish that every intermediate claim has been independently checked.

Imported mathematical results retain their attribution and assumptions. A classical theorem used as an explicit import is an ordinary mathematical dependency. For compatibility with Phase II, a \emph{source dependency} means an unresolved reliance on a particular extracted source assertion or certificate contract, rather than every imported theorem. This convention prevents a label intended to flag an unreviewed research input from also flagging routine uses of established linear algebra.

Novelty fields distinguish ``known in the collection,'' ``new to the selected record,'' and ``literature novelty unresolved.'' A search that fails to locate a prior statement does not prove novelty. Significance is assessed separately in \cref{sec:significance}. The case study attributes inherited lemmas and classical inputs separately from the certificate construction and transfer developed during the run.

\subsection{Conditional support and safe graph updates}

Consider a proposed route
\[
 A\Longrightarrow B,\qquad B\land C\Longrightarrow D.
\]
The route reaches $D$ only after $C$ is established and the two occurrences of $B$ have compatible meanings and scopes. If $C$ is merely conjectured, the graph may record a conditional argument for $D$ but cannot silently mark $D$ as proved. A counterexample to one proposed $C$ eliminates that route in its stated scope, not every possible proof of $D$.

An update therefore has two effects. It appends historical records and changes the current view of supported routes. Historical storage can be append-only while current admissibility changes nonmonotonically. If a premise is withdrawn, affected descendants are marked for reassessment. A conclusion with an alternative valid derivation may remain supported. The system must assess individual derivations rather than indiscriminately invalidate every connected node.

For ordinary finite arguments, the admitted proof dependencies should be acyclic. Exploration can revisit a target many times, and the historical graph can contain feedback patterns. Such feedback cannot establish two unsupported claims by having them cite each other. Mathematical induction or a justified fixed-point argument is represented through its explicit theorem and discharge conditions, not through an unexplained cycle of assertions.

Structural validators can detect unresolved identifiers, missing evidence files, malformed scopes, or circular dependencies. They cannot generally decide whether an informal inference is valid or whether two hypotheses are equivalent. The graph is a means of exposing these questions to tools and reviewers; its structural validity is a necessary bookkeeping property, not a proof certificate for the mathematics.

\subsection{The research target as a first-class record}

Each target specifies a statement, its parameter regime, the objective-relevant gap it addresses, candidate source arguments, an acceptance criterion, and the consequences that success would enable. It also records an investigation budget and the evidence available when the target was selected. An anticipated downstream implication must itself be justified or listed as conditional.

For example, the target in the case study asks for matching upper and lower bounds on the degree of a polynomial identity in the original tensor-decomposition parameters. A proof using inverses or normalized variables must account for the cost of returning to that parameter model. The target specifies an infinite family, so a certificate for one small tensor supplies a new input rather than completing the objective. Recording these requirements makes the next useful question explicit before significance is assigned.

The target frontier need not contain only theorem statements. It can include a request to test a proposed mechanism against a known obstruction, compare two quantifier orders, formalize a critical lemma, or reproduce a computational certificate. These activities have research value when their consequences for the objective are explicit.
